% Generated draft: requires mathematical and source review.
\documentclass[11pt]{article}
\usepackage[margin=25mm]{geometry}
\usepackage{fontspec}
\setmainfont{Latin Modern Roman}
\usepackage{amsmath,amssymb,mathtools,unicode-math}
\setmathfont{Latin Modern Math}
\usepackage{enumitem,needspace,xurl,hyperref}
\setlength{\parindent}{0pt}
\setlength{\parskip}{5pt}
\setlength{\emergencystretch}{3em}
\newcommand{\N}{\mathbb{N}}
\newcommand{\Z}{\mathbb{Z}}
\newcommand{\Q}{\mathbb{Q}}
\newcommand{\R}{\mathbb{R}}
\newcommand{\C}{\mathbb{C}}
\newcommand{\F}{\mathbb{F}}
\newcommand{\A}{\mathbb{A}}
\newcommand{\PP}{\mathbb{P}}
\newcommand{\E}{\mathbb{E}}
\newcommand{\eps}{\varepsilon}
\newcommand{\dd}{\,\mathrm{d}}
\begin{document}
\section*{UnsolvedMath ID: 800174}
\textbf{Source problem code:} catalogue entry
\subsection{Definitions and mathematical statement}
Fix a prime \(p>2\) and a power \(q\) of \(p\). Let \(\psi\) denote the additive character occurring in the source, and let \(\mathcal L(u,f,\psi)\) denote the associated Artin--Schreier \(L\)-function. The source does not specify the domain or nontriviality of \(\psi\), or the normalization of \(\mathcal L\).

For positive integers \(d\), define the following families, interpreting the misplaced braces in the supplied transcription as restrictions on \(d\):
\[
\mathcal{AS}^{0}_d=\{f\in\F_q[x]:\deg f=d\},\qquad \gcd(d,p)=1,
\]
\[
\mathcal{AS}^{0,\mathrm{odd}}_d
=\{f\in\F_q[x]:\deg f=d,\ f(x)=-f(-x)\},
\qquad \gcd(d,2p)=1,
\]
and
\[
\begin{split}
\mathcal{AS}^{\mathrm{ord}}_d
=\bigl\{f\in\F_q(x):\;&f=h/g,\quad h,g\in\F_q[x],\quad g\ne0,\\
&\gcd(h,g)=1,\quad g\text{ squarefree},\\
&\max(\deg h,\deg g)=d,\quad \deg g\in\{d,d-1\}\bigr\}.
\end{split}
\]
Here the condition called \(\deg f=d\) in the source is explicitly \(\max(\deg h,\deg g)=d\) for a reduced representation.

For each family \(\mathcal F_d\), retain the source's moment notation in the family-indexed form
\[
M_{k,\mathcal F}(d)
=\left\langle\mathcal L(q^{1/2},f,\psi)^k\right\rangle_{f\in\mathcal F_d}.
\]
The brackets denote the source's family average; its measure and normalization are not specified. The argument \(q^{1/2}\) is reproduced literally. Whether this is the intended evaluation point depends on the unspecified convention for \(\mathcal L\) and requires clarification.

\textbf{Moment targets.} With \(q\) fixed and \(d\to\infty\), subject to the displayed restrictions on \(d\), compute the moments for the following orders:
\[
\begin{array}{c|c}
\mathcal F_d&k\\ \hline
\mathcal{AS}^{\mathrm{ord}}_d&1,2,3,4,5,6\\
\mathcal{AS}^{0}_d&1,2,3,4\\
\mathcal{AS}^{0,\mathrm{odd}}_d&1,2.
\end{array}
\]
The source does not specify whether ``compute'' requires an exact formula, a leading asymptotic, further terms, or a particular error bound.

\textbf{One-level density targets.} In the same fixed-\(q\), growing-\(d\) regime, compute the family-averaged one-level density of zeros against a Schwartz test function \(\phi\), with the following Fourier-support extensions:
\begin{enumerate}
\item For \(\mathcal{AS}^{\mathrm{ord}}_d\), the source lists
\[
\operatorname{supp}\widehat\phi\subset(-1-\delta,1+\delta)
\quad\text{for some }\delta>0,
\]
or
\[
\operatorname{supp}\widehat\phi\subset(-2-\delta,2+\delta)
\quad\text{for some }\delta>0.
\]
The source does not clarify whether these are alternative milestones or separately required targets.
\item For \(\mathcal{AS}^{0}_d\), the listed support condition is
\[
\operatorname{supp}\widehat\phi
\subset\bigl(-(2-2/p)-\delta,\ 2-2/p+\delta\bigr).
\]
\item For \(\mathcal{AS}^{0,\mathrm{odd}}_d\), the listed support condition is
\[
\operatorname{supp}\widehat\phi
\subset\bigl(-(1-1/p)-\delta,\ 1-1/p+\delta\bigr).
\]
\end{enumerate}
In the last two items, the supplied text does not explicitly quantify \(\delta\); interpreting it as a positive support extension requires confirmation. The source also leaves the zero scaling, treatment of periodic zero ordinates, Fourier-transform convention, averaging measure, and precise density asymptotic unspecified. Thus these are faithfully recorded research targets, rather than fully normalized density assertions.
\subsection{Short English statement}
For fixed \(q\) of odd characteristic, determine the listed moments and extend the Fourier-support range for one-level zero-density results in the ordinary, polynomial, and odd-polynomial Artin--Schreier families as \(d\to\infty\). The supplied formulation leaves essential normalization and quantifier choices unresolved, including the literal moment evaluation point \(q^{1/2}\). The registry labels the problem open; this is supplied status information, not an independently verified assessment.
\subsection{Sources}
\begin{itemize}
\item Supplied local UnsolvedMath registry record, permanent ID \(800174\), exact source problem code \texttt{catalogue entry}, in the AMR Open Problem Lists collection. The record gives status \texttt{open}, proposed year \(2024\), and \texttt{proposed\_by}: Hamieh, Kadiri, Ng. Its title is ``Moments of L-functions --- Moments of the Riemann zeta function --- Problem 25,'' although its mathematical text concerns Artin--Schreier \(L\)-functions.
\item Supplied source-list attribution: Hamieh, Kadiri, Ng, \emph{Moments of L-functions Problem List} (2024), Problem 25, \url{https://arxiv.org/abs/2405.17800}. The supplied background records access on \(2026\text{-}10\text{-}16\) and extraction method \texttt{arxiv\_latex}; these are supplied provenance statements, not evidence of independent source access for this draft.
\item The supplied problem text ends with the attribution ``Alexei Entin, Tel Aviv University.'' This attribution is retained separately from the registry's attribution to Hamieh, Kadiri, Ng.
\end{itemize}
\end{document}
